\documentclass[11pt]{article}
\usepackage[a4paper,margin=1in]{geometry}
\usepackage[T1]{fontenc}
\usepackage{lmodern,microtype}
\usepackage{amsmath,amssymb,amsthm,mathtools}
\usepackage{booktabs,enumitem}
\usepackage{xcolor}
\usepackage[numbers,sort&compress]{natbib}
\usepackage[colorlinks=true,linkcolor=blue!55!black,citecolor=blue!55!black]{hyperref}

\newtheorem{theorem}{Theorem}[section]
\newtheorem{proposition}[theorem]{Proposition}
\newtheorem{corollary}[theorem]{Corollary}
\newtheorem{lemma}[theorem]{Lemma}
\theoremstyle{definition}
\newtheorem{definition}[theorem]{Definition}
\theoremstyle{remark}
\newtheorem{remark}[theorem]{Remark}

\title{Deterministic-Numerator Analytic-Tail Certificate}
\author{Bin Wang}
\date{July 2026}

\begin{document}
\maketitle

\begin{abstract}
RH-251 left the deterministic numerator tail as a separate obligation in the
finite-head/analytic-tail interface.  RH-46 supplies a stronger analytic input
than the finite coefficient dictionary alone: the deterministic one-step
numerator $G$ is holomorphic and nonzero on $|z|<\lambda$, where
$\lambda=1.6785735104283177\ldots$.  After Hardy scaling by $r_H=0.85$, the
logarithm of $G(z/r_H)$ therefore has radius at least
$\rho_*=r_H\lambda=1.42678748386407\ldots$.  We prove an all-order Cauchy
tail bound for the target coefficients and its exponential determinant
conversion.  The boundary supremum entering the bound is not numerically
certified, so no cloud or quotient uniformity is claimed.
\end{abstract}

\section{Input and scope}

Write the deterministic factor from RH-46 as
\begin{equation}
 \widehat D_{0,\mathrm{bulk},2}(z)=\frac{G(z)}{1-z^2/\lambda},
 \qquad G(0)=1,
 \label{eq:factor}
\end{equation}
where $G$ is holomorphic and nonzero on $|z|<\lambda$ \citep{WangRH46}.
RH-243 identifies the one-step Hardy-scaled trace-style target by
\begin{equation}
 \log G_H(z)=-\sum_{n\ge2}\frac{a_n}{n}z^n,
 \qquad G_H(z):=G(z/r_H),
 \qquad r_H=0.85.
 \label{eq:target}
\end{equation}
The coefficients $a_n$ are deterministic target coefficients, not yet the
coefficients of a selected noisy cloud \citep{WangRH243}.  We set $a_1:=0$,
so coefficient estimates below may be stated uniformly for $n\ge1$.

\begin{proposition}[Scaled zero-free disk]
\label{prop:radius}
The function $G_H$ is holomorphic and nonzero on
\begin{equation}
 |z|<\rho_*:=r_H\lambda=1.42678748386407\ldots>1.
 \label{eq:rho}
\end{equation}
Consequently there is a unique holomorphic logarithm $L_H$ on this disk with
$L_H(0)=0$, and $L_H=\log G_H$ in \eqref{eq:target}.
\end{proposition}

\begin{proof}
The substitution $z\mapsto z/r_H$ maps $|z|<r_H\lambda$ into
$|z|<\lambda$.  Holomorphicity and nonvanishing are preserved.  The disk is
simply connected and $G_H(0)=1$, so the normalized logarithm exists and is
unique.
\end{proof}

\section{The all-order tail theorem}

\begin{definition}[Boundary budget]
For $0\le R<S<\rho_*$ define
\begin{equation}
 M_S:=\max_{|z|=S}|L_H(z)|,
 \qquad
 \Phi_N(R,S):=\frac{(R/S)^N}{1-R/S},
 \quad N\ge1.
 \label{eq:budget}
\end{equation}
The first omitted order is $N$; thus the tail starts at $n=N$.
\end{definition}

\begin{theorem}[Analytic deterministic-target tail]
\label{thm:tail}
For every $0\le R<S<\rho_*$ and every integer $N\ge1$,
\begin{equation}
 \frac{|a_n|}{n}\le M_S S^{-n}\quad(n\ge1),
 \qquad
 \sum_{n\ge N}\frac{|a_n|R^n}{n}
 \le M_S\Phi_N(R,S).
 \label{eq:tail}
\end{equation}
In particular, since $\rho_*>1$, the all-order deterministic target tail on
the unit disk exists for every $1<S<\rho_*$.
\end{theorem}

\begin{proof}
By Cauchy's coefficient estimate applied to the holomorphic function $L_H$ on
$|z|<S$, the coefficient of $z^n$ in $L_H$ has modulus at most
$M_SS^{-n}$.  Equation \eqref{eq:target} says that this coefficient is
$-a_n/n$.  Multiplying by $R^n$ and summing the geometric series proves
\eqref{eq:tail}.
\end{proof}

\begin{corollary}[Determinant conversion]
\label{cor:exp}
If $E_N(R,S):=M_S\Phi_N(R,S)$, then truncating the target logarithm after
order $N-1$ incurs relative multiplicative error at most
\begin{equation}
 \left|\exp\!\left(-\sum_{n\ge N}a_nz^n/n\right)-1\right|
 \le e^{E_N(R,S)}-1,
 \qquad |z|\le R.
 \label{eq:exp-error}
\end{equation}
\end{corollary}

\begin{proof}
Theorem~\ref{thm:tail} bounds the modulus of the omitted logarithm by
$E_N(R,S)$.  The elementary inequality $|e^w-1|\le e^{|w|}-1$ gives the
claim.
\end{proof}

\begin{remark}[What the theorem does not provide]
The theorem is an exact analytic interface, but $M_S$ is an unknown boundary
constant.  A finite evaluation of the order-$2$--$12$ polynomial cannot bound
the full $L_H$ on the same circle.  Thus the result is not a numerical
uniform tail certificate for the moving cloud and does not invoke RH-240.
\end{remark}

\section{Finite diagnostic}

We read the archived RH-243 coefficients through order twelve and evaluate
the truncated logarithm on $4096$ equally spaced points for
$S\in\{1.05,1.15,1.25,1.35\}$.  These values are diagnostics only; the exact
quantity $M_S$ remains unbounded by the experiment.  The geometric factor for
the unit-disk tail beginning at order $13$ is shown below.

\begin{center}
\small
\begin{tabular}{rrrr}
\toprule
$S$ & truncated $\max|L_H|$ & $\Phi_{13}(1,S)$ & status\\
\midrule
1.05 & 0.579287 & 11.136748 & diagnostic only\\
1.15 & 0.797017 & 1.246048 & diagnostic only\\
1.25 & 1.106719 & 0.274878 & diagnostic only\\
1.35 & 1.559934 & 0.077970 & diagnostic only\\
\bottomrule
\end{tabular}
\end{center}

The last column is essential: multiplying the displayed truncated supremum by
$\Phi_{13}$ is not a proof of the tail bound.  It only illustrates the radius
budget once a rigorous $M_S$ is supplied.

\section{Gate and route audit}

The new exact statement closes the deterministic-target analytic-tail
existence obligation on every strict subdisk of radius $\rho_*$.  It does not
identify the current noisy cloud with $a_n$, does not supply a uniform quotient
block estimate, and does not produce a finite numerical value for $M_S$.
Therefore Gates A--E remain false/open.  No Hilbert--Polya operator, Riemann
zero identification, zeta-divisor equality, or RH implication is claimed.

The next result should either interval-certify $M_S$ on a useful circle or
provide a genuinely expanded resolved candidate window.  Reweighting the
frozen RH-248 shell class would not be a new input.

\bibliographystyle{plainnat}
\bibliography{references}
\end{document}
